\chapter*{Tóm tắt đề tài}
\addcontentsline{toc}{chapter}{Tóm tắt đề tài}
\markboth{TÓM TẮT ĐỀ TÀI}{}

Thị trường không gian làm việc chia sẻ (Co-working Space) tại Việt Nam đang tăng trưởng nhanh, nhưng phần lớn các đơn vị vận hành vẫn quản lý thủ công qua bảng tính và tin nhắn, dẫn đến các vấn đề: xung đột trùng lịch đặt chỗ, sai sót đối soát thanh toán, chính sách hủy thiếu nhất quán và bỏ ngỏ giá trị kết nối cộng đồng giữa các thành viên.

Đề tài \textbf{``Phát triển hệ thống quản lý văn phòng chia sẻ (CoSpace)''} xây dựng một nền tảng số hóa toàn diện cho chuỗi Co-working Space đa chi nhánh, phục vụ bốn nhóm tác nhân: Khách hàng, Nhân viên quầy, Quản lý chi nhánh và Quản trị viên hệ thống. Hệ thống bao gồm: sơ đồ mặt bằng tương tác SVG, đặt chỗ thời gian thực, tích hợp thanh toán VietQR PayOS và MoMo, chính sách hủy--hoàn tiền tự động theo bậc thời gian, check-in/check-out bằng mã QR, gợi ý đối tác kinh doanh dựa trên thuật toán Jaccard, và trợ lý ảo Google Gemini AI có rào chắn xác nhận giao dịch.

Về công nghệ, hệ thống sử dụng Spring Boot 3.1 (Java 17), React 18 với TypeScript và Tailwind CSS, PostgreSQL trên Supabase, di trú lược đồ bằng Flyway. Bài toán tương tranh được giải quyết bằng cơ chế hai lớp: khóa tư vấn \texttt{pg\_advisory\_xact\_lock} tại tầng ứng dụng kết hợp ràng buộc loại trừ \texttt{EXCLUDE USING gist} tại tầng cơ sở dữ liệu. Webhook thanh toán được thiết kế bất biến (Idempotent) qua trạng thái giao dịch và khóa dòng bi quan.

Sản phẩm đạt quy mô: \textbf{30 bảng CSDL}, \textbf{26 lớp điều khiển} với \textbf{133 API endpoints}, \textbf{33 trang giao diện}, \textbf{307 ca kiểm thử tự động} Back-end (0 lỗi). Hệ thống được đóng gói Docker đa giai đoạn và triển khai trên Render, Vercel và Supabase.

\vspace{0.3cm}
\noindent Báo cáo gồm 8 chương:
\begin{itemize}
    \item \textbf{Chương 1:} Tổng quan -- bối cảnh, bài toán, mục tiêu, phạm vi và kế hoạch thực hiện.
    \item \textbf{Chương 2:} Khảo sát hiện trạng -- phân tích mô hình Co-working và các hệ thống liên quan.
    \item \textbf{Chương 3:} Cơ sở lý thuyết -- kiến trúc 3 tầng, RESTful, JWT, Advisory Locks, Jaccard, Gemini.
    \item \textbf{Chương 4:} Phân tích yêu cầu -- 4 vai trò, 20 yêu cầu chức năng, yêu cầu phi chức năng, Use-case.
    \item \textbf{Chương 5:} Thiết kế hệ thống -- Class Diagram, Sequence Diagram, ERD, State Machine.
    \item \textbf{Chương 6:} Hiện thực và Giao diện -- Service Specs, thuật toán tương tranh, Docker, 12 màn hình.
    \item \textbf{Chương 7:} Kiểm thử và đánh giá -- 307 ca kiểm thử tự động, Race conditions, UAT.
    \item \textbf{Chương 8:} Tổng kết -- kết quả, hạn chế và hướng phát triển.
\end{itemize}

\vspace{0.4cm}
\noindent\textbf{Từ khóa:} Co-working Space, đặt chỗ trực tuyến, kiểm soát tương tranh, Advisory Lock, EXCLUDE USING gist, VietQR, Webhook Idempotent, Jaccard, Gemini AI, Spring Boot, React.
